\section{Motor de datos: el ciclo cerrado de datos que va del fracaso al éxito}

Xie Chen (谢晨) dice que los datos más eficaces son los de «primero fallar y luego tener éxito». Esta afirmación es importante. Las trayectorias exitosas le indican al modelo qué debe hacer; las trayectorias que fallan y luego se corrigen le indican por qué se equivocó, cómo recuperarse y cómo aprender del error. Tanto la conducción autónoma como los robots necesitan este ciclo cerrado: despliegue, recolección, minería, entrenamiento, evaluación y nuevo despliegue.

\begin{figure}[H]
\centering
\includegraphics[width=0.94\textwidth]{data-engine-loop.png}
\caption{Ciclo cerrado de datos / Data Engine: despliegue, recolección, minería, entrenamiento, evaluación, simulación.}
\end{figure}

\begin{knowledgebox}{Lectura de la figura: lo esencial del Data Engine es la continuidad}
En la figura, cada ronda de despliegue trae de vuelta retroalimentación real; los casos de fracaso se extraen y se convierten en datos; después del entrenamiento se validan mediante evaluación en simulación y, por último, se vuelve a desplegar. La conclusión que sustenta es esta: los datos no son una compra única, sino el resultado de que el sistema opere de forma continua.
\end{knowledgebox}

\subsection{Lo que enseña el motor de datos de Tesla}

El motor de datos de Tesla se alimenta de una gran cantidad de vehículos en el lado del dispositivo. Los autos operan en el mundo real, generan escenarios de cola larga y retroalimentación de conducción, que regresan a la nube para el entrenamiento y actualizan las capacidades en el lado del dispositivo. Esta lógica no se cumple del todo en los robots, porque la escala de robots en el lado del dispositivo es insuficiente. Por ello, la industria robótica debe buscar un nuevo motor de datos: una combinación de simulación, síntesis, teleoperación, minería de casos de fracaso, benchmark y despliegue real a pequeña escala.

\begin{importantbox}{La esencia del motor de datos}
Un motor de datos no consiste en «tener muchos datos», sino en contar con un ciclo capaz de detectar errores de manera continua, generar datos complementarios, validar las mejoras y volver a desplegar.
\end{importantbox}

\subsection{Resumen de la sección}

Los datos de alta calidad provienen de un ciclo cerrado. Tanto los robots como los LLM necesitan retroalimentación, pero la retroalimentación de los robots es más costosa y depende más del entorno y de la evaluación.
